\documentclass[12pt]{article}
\usepackage[utf8]{inputenc}
\usepackage[T1]{fontenc}
\usepackage[margin=1in]{geometry}
\usepackage{setspace}
\usepackage{xcolor}
\usepackage[colorlinks=true, urlcolor=blue!60!black, linkcolor=blue!60!black]{hyperref}
\setlength{\parindent}{2em}
\setlength{\parskip}{1.25em}
\renewcommand{\baselinestretch}{1.0}
\begin{document}
\begin{center}
{\LARGE \textit{United States v. Comstock}, 560 U.S. 126 (2010)\par}
\vspace{0.8em}
{\large Case Brief --- Four-Page Edition\par}
\end{center}
\vspace{1em}

\section{Memory Jogger}

Federal prisoners nearing release challenged a statute allowing their indefinite civil commitment as sexually dangerous; the Court held 7--2 that the Necessary and Proper Clause authorized it as a rational, narrow incident of Congress's power to run federal prisons and act as a responsible custodian (\href{https://supreme.justia.com/cases/federal/us/560/126/}{560 U.S. 126}).

\section{Detailed Case Facts}

In 2006, as part of the Adam Walsh Child Protection and Safety Act, Congress enacted 18 U.S.C. \S~4248, authorizing a federal court to order indefinite civil commitment of a ``sexually dangerous'' person already in Bureau of Prisons custody---including prisoners nearing the end of their sentences, defendants found incompetent to stand trial, and persons whose charges were dismissed because of their mental condition. The government must prove by clear and convincing evidence that the person has engaged or attempted to engage in sexually violent conduct or child molestation, suffers from a serious mental illness, abnormality, or disorder, and as a result would have serious difficulty refraining from such conduct if released. The Attorney General must make all reasonable efforts to transfer custody to the person's State of domicile or trial; federal confinement continues only if no State assumes responsibility, and only until the person is no longer dangerous.

\section{Procedural History}

The government sought commitment of Graydon Earl Comstock, Jr., and four other detainees in the Eastern District of North Carolina. The District Court held that Congress lacked Article I power and dismissed. The Fourth Circuit affirmed on that ground alone, not reaching due process or equal protection. The Supreme Court granted certiorari limited to Article I, heard argument on January 12, 2010, and decided the case on May 17, 2010.

\section{Judicial Votes}

Reversed and remanded, \textbf{7--2}.
\begin{itemize}
\item \textbf{Majority:} Justice Breyer, joined by Chief Justice Roberts and Justices Stevens, Ginsburg, and Sotomayor.
\item \textbf{Concurring in the judgment:} Justice Kennedy; Justice Alito.
\item \textbf{Dissent:} Justice Thomas, joined by Justice Scalia in all but Part III-A-1-b.
\end{itemize}

\section{Holding}

The Necessary and Proper Clause empowers Congress to enact \S~4248 as ``a means that is rationally related to the implementation of a constitutionally enumerated power.'' Relying on five considerations taken together, the Court reversed and remanded for consideration of the preserved due process and equal protection claims.

\section{Analysis of Opinions}

\textbf{Majority (Breyer).} \textit{First}, under \textit{McCulloch} and \textit{Sabri}, the Clause authorizes laws ``convenient, or useful'' or ``conducive'' to exercising enumerated powers; Congress may, by implication, create federal crimes, build prisons, and regulate prison safety. \textit{Second}, Congress has provided care and commitment for mentally ill federal prisoners since 1855, including the 1949 statute upheld in \textit{Greenwood}, so \S~4248 was a modest addition. \textit{Third}, a custodian that lawfully holds a prisoner acts responsibly in preventing the release of a dangerous one. \textit{Fourth}, powers delegated by the Clause are not ``reserved to the States,'' and \S~4248(d) defers to States willing to take custody. \textit{Fifth}, rejecting a rule against acting more than one implied step from an enumerated power, Breyer found the links not too attenuated because the statute reaches only a small class already in federal custody and confers no general police power.

\textbf{Kennedy (concurring in the judgment).} The statute is a narrow exercise of custodial responsibility that respects state primacy, but importing ``rational basis'' rhetoric from due process review obscures federalism's limits; the strength of the link to an enumerated power is what matters.

\textbf{Alito (concurring in the judgment).} ``Necessary'' requires a ``substantial link,'' satisfied here because federal imprisonment often severs inmates' ties to any State, creating a gap the federal government helped produce.

\textbf{Thomas (dissenting).} The statute is neither ``plainly adapted'' to an identifiable enumerated power nor consistent with the Constitution's spirit; once a sentence ends, the power behind the conviction is exhausted, and committing dangerous citizens is a local police power.

\textbf{Later developments.} In \textit{NFIB v. Sebelius} (2012), Chief Justice Roberts distinguished \textit{Comstock} as exercising authority ``derivative of, and in service to, a granted power'' and found the individual mandate not ``proper.'' In \textit{United States v. Kebodeaux} (2013), the Court applied \textit{Comstock} to uphold registration duties for a former serviceman subject to a preexisting federal obligation.

\section{Examples: Future Applications}

\begin{enumerate}
\item \textbf{Same-Side (Valid): Violently Dangerous Federal Prisoners.} Extending commitment to mentally ill prisoners who would pose a serious danger of violence, with the same proof and state-transfer provisions, satisfies every \textit{Comstock} consideration.
\item \textbf{Same-Side (Valid): Contagious Prisoner.} Holding a highly contagious inmate in isolation past his release date until he is no longer contagious, unless his State takes custody, is a narrow custodial measure.
\item \textbf{Opposite-Side (Invalid): State Felons via a Firearm Nexus.} Committing any past state felon who once bought a gun that moved in interstate commerce, though never in federal custody, lacks a custodial link and intrudes on state police power.
\item \textbf{Opposite-Side (Invalid): People Never Charged.} Committing people the FBI investigated but never charged reaches beyond federal custody, the feature that kept \S~4248's links from being too attenuated.
\item \textbf{Fence-Sitter: Registration for Long-Released Offenders.} Requiring registration from federal sex offenders released decades ago without conditions lacks \textit{Kebodeaux}'s continuing obligation but is less intrusive than commitment.
\end{enumerate}

\section{Critique}

\textbf{The stopping-point problem.} Breyer's five-factor ``taken together'' test is pragmatic and flexible, but critics echo Thomas: if Congress may stack implied power on implied power, from enumerated power to federal crime, prison, prison medical care, and post-sentence commitment, what principled stopping point remains?

\textbf{``Proper'' as a federalism check.} Kennedy's and Alito's concurrences, and Roberts's later opinion in \textit{NFIB}, recast ``proper'' as an independent check preventing Congress from converting incidental custodial powers into a general police power.

\section{Key Quotations}

\begin{quote}
``Taken together, these considerations lead us to conclude that the statute is a `necessary and proper' means of exercising the federal authority that permits Congress to create federal criminal laws, to punish their violation, to imprison violators, to provide appropriately for those imprisoned, and to maintain the security of those who are not imprisoned but who may be affected by the federal imprisonment of others.''
\end{quote}
\noindent --- Justice Breyer, Opinion of the Court, 560 U.S. at 149.

\begin{quote}
``The Necessary and Proper Clause does not give Congress carte blanche. Although the term `necessary' does not mean `absolutely necessary' or indispensable, the term requires an `appropriate' link between a power conferred by the Constitution and the law enacted by Congress. ... Here, there is a substantial link to Congress' constitutional powers.''
\end{quote}
\noindent --- Justice Alito, concurring in the judgment, 560 U.S. at 158.

\begin{quote}
``Concluding that a relation can be put into a verbal formulation that fits somewhere along a causal chain of federal powers is merely the beginning, not the end, of the constitutional inquiry. ... The terms `rationally related' and `rational basis' must be employed with care, particularly if either is to be used as a stand-alone test.''
\end{quote}
\noindent --- Justice Kennedy, concurring in the judgment, 560 U.S. at 150--51.

\end{document}
